\chapter{Hand Washing}
\index{Hand Washing}

\section*{Definition and Overview}
Hand washing is the simple act of cleaning the hands with soap and water, or with an alcohol-based hand sanitiser, to remove germs, dirt, and chemicals. It is one of the most effective ways to prevent the spread of infections, including colds, influenza, COVID-19, gastroenteritis, and foodborne illness. Correct hand washing involves wetting the hands, applying soap, scrubbing all surfaces for at least 20 seconds, rinsing, and drying. Hands should be washed at key times: after using the toilet, before and after preparing food, before eating, after coughing or sneezing, after touching animals, and when caring for someone who is unwell.

\section*{Epidemiology}
Hand hygiene has a major impact on public health. Around 80% of common infections are transmitted by hands, and studies show that regular hand washing reduces the risk of respiratory infections by around 16--21% and diarrhoeal illness by up to 40%. In health care settings, hand hygiene reduces hospital-acquired infections substantially, and local hospitals audit and report hand hygiene compliance. During the COVID-19 pandemic, hand hygiene was promoted worldwide as a key preventive measure. Despite its effectiveness, hand washing rates remain suboptimal in many settings.

\section*{Aetiology and Pathophysiology}
Hand washing works by physically removing germs from the skin.

\begin{itemize}
    \item \textbf{Germs on hands:} hands pick up bacteria, viruses, and other germs from surfaces, people, animals, and body fluids
    \item \textbf{Mechanism:} soap and friction lift and remove germs, dirt, and oils; water rinses them away
    \item \textbf{Soap:} dissolves the fatty envelope of many viruses and helps detach bacteria
    \item \textbf{Alcohol sanitiser:} at least 60% alcohol rapidly inactivates most germs when hands are not visibly dirty
    \item \textbf{Transmission:} unwashed hands spread germs to the mouth, eyes, and nose, to food, and to other people
    \item \textbf{Key sites:} fingertips and under the nails harbour most germs, so thorough scrubbing matters
    \item \textbf{Timing:} at least 20 seconds of scrubbing is recommended
\end{itemize}

\section*{Clinical Features}
Hand washing is relevant to preventing many illnesses. Situations that require hand washing include:

\begin{itemize}
    \item After using the toilet, changing nappies, or helping a child use the toilet
    \item Before, during, and after preparing food, especially handling raw meat
    \item Before eating
    \item After coughing, sneezing, or blowing the nose
    \item After touching animals, animal feed, or animal waste
    \item After handling rubbish or waste
    \item When caring for someone who is sick
    \item After treating a cut or wound
    \item When hands are visibly dirty or greasy
\end{itemize}

\section*{Differential Diagnosis}
Hand hygiene matters most in situations where infection risk is highest.

\begin{itemize}
    \item Food handling and preparation, to prevent foodborne illness
    \item Health care and caring for vulnerable people
    \item Childcare settings, where infections spread rapidly
    \item After contact with high-risk surfaces, such as public transport and gym equipment
    \item During outbreaks of respiratory and gastrointestinal infections
    \item After handling pets, reptiles, and raw pet food
    \item For people with weakened immunity, who are at higher risk
\end{itemize}

\section*{When to Seek Medical Care}
Hand washing itself does not require medical care. See a doctor if you develop symptoms of an infection that might have been preventable, such as persistent diarrhoea, or if caring for someone unwell and you have health concerns.

\textbf{Contact your local emergency services for:}
\begin{itemize}
    \item This topic concerns prevention rather than emergencies; however, seek urgent care for any severe or rapidly worsening illness
    \item Severe dehydration from gastroenteritis
    \item Signs of sepsis, including high fever, confusion, or collapse
\end{itemize}

\section*{Diagnosis and Investigations}
No tests are involved in hand washing; the focus is on correct technique.

\begin{itemize}
    \item \textbf{Correct technique:} wet hands, apply soap, scrub all surfaces (backs of hands, between fingers, under nails) for 20 seconds, rinse, and dry
    \item \textbf{Hand sanitiser technique:} apply enough to cover all surfaces and rub until dry
    \item \textbf{Choosing the product:} soap and water are preferred when hands are visibly dirty; sanitiser is a good alternative otherwise
    \item \textbf{Drying:} drying matters, as germs transfer more easily from wet hands
    \item \textbf{In health care:} specific hand hygiene protocols with audits
\end{itemize}

\section*{Management}
Making hand washing a habit is the most effective management of infection risk.

\begin{itemize}
    \item \textbf{Wash at key times:} after the toilet, before food, after coughing or sneezing, and after contact with animals or sick people
    \item \textbf{Use soap and water when possible:} especially when hands are visibly dirty, after using the toilet, and after handling raw meat
    \item \textbf{Use sanitiser on the go:} keep alcohol-based sanitiser (at least 60% alcohol) in bags, cars, and workplaces
    \item \textbf{Teach children:} model and teach hand washing habits early
    \item \textbf{In food preparation:} wash hands before, during, and after handling food
    \item \textbf{In care settings:} follow health care hand hygiene protocols
    \item \textbf{Keep nails short:} reduces the harbour for germs
    \item \textbf{Use paper towels or a clean towel:} for drying
\end{itemize}

\section*{Complications}
\begin{itemize}
    \item Skin irritation and dryness from frequent washing, managed with moisturiser
    \item Poor technique, such as brief rinsing without soap, which is much less effective
    \item Missed key times, which leaves the highest-risk contacts unprotected
    \item Reliance on sanitiser when hands are visibly dirty
    \item False reassurance if washing is not done at the right times
    \item Infection transmission when hands are not washed
\end{itemize}

\section*{Prognosis and Outcomes}
Consistent hand washing produces large, measurable reductions in infection rates. Communities and households that practise good hand hygiene have fewer respiratory infections, fewer gastrointestinal illnesses, and less spread of infection to vulnerable people. In health care, high hand hygiene compliance is associated with fewer hospital-acquired infections. Hand washing is one of the cheapest, simplest, and most effective health measures available, and its benefits are lifelong.

\section*{Prevention and Self-Care}
\begin{itemize}
    \item Wash hands for at least 20 seconds with soap and water
    \item Wash after using the toilet and before eating or preparing food
    \item Wash after coughing, sneezing, and touching animals
    \item Use alcohol-based sanitiser when soap and water are unavailable
    \item Teach children to wash their hands properly
    \item Use a moisturiser to prevent dry skin from frequent washing
    \item Avoid touching your face with unwashed hands
    \item Wash hands before and after caring for someone who is unwell
    \item Keep hand sanitiser available at home, work, and in the car
\end{itemize}

\section*{Sources and Further Reading}
The following reputable sources provide further information for healthcare students and patients seeking reliable, current advice about hand washing.

\begin{itemize}
    \item Australian Government Department of Health and Aged Care. \textit{Hand hygiene.} https://www.health.gov.au/
    \item Australian Commission on Safety and Quality in Health Care. \textit{Hand hygiene in health care.} https://www.safetyandquality.gov.au/
    \item Healthdirect Australia. \textit{Hand washing: how to do it properly.} https://www.healthdirect.gov.au/
    \item Food Safety Information Council. \textit{Hand washing and food safety.} https://foodsafety.asn.au/
    \item Centers for Disease Control and Prevention. \textit{Handwashing: clean hands save lives.} https://www.cdc.gov/
    \item World Health Organization. \textit{Hand hygiene.} https://www.who.int/
\end{itemize}
